% !TEX encoding = UTF-8 Unicode
% !TEX root = ../curriculum.tex

\cvsection{Experiencias: Detalle de contribuciones}

\begin{cvparagraph}
    Diseñé e implemente un sistema de carga de datos modular que permitía integrar a
    una BBDD PostgreSQL, Buckets S3 y Google Drive; datos desde planillas de cálculo
    de Office, planillas en la nube de Google Sheets, LibreOffice Calc y GMail. La
    arquitectura es lo suficientemente flexible para que cada integración sea
    completamente independiente y testeable, confiriendo mucha flexibilidad al uso y
    desarrollo.

    Como frontend, por simplicidad y constrains de tiempo se utilizaron notebooks
    jupyter, pero la API permitía integrarlo a cualquier librería de escritorio tipo
    TK, Qt, Gtk, etc. o a cualquier framework web.

    Para el deploy de esta aplicación, levanté una pequeña pipeline en Bitbucket que
    aplicaba formateo automático (Black), ejecutaba los tests en una instancia que
    también se generaba automáticamente y en caso de aprobar, generaba un artefacto
    que utilizaba el equipo de analítica.

    El sistema era tan versátil que incluso se utilizó para carga pseudo manual de
    datos a través de correos-tipo que podían enviar los clientes.

    En base a este desarrollo pasamos de integrar de 2 a 3 componentes semanales a
    15 diarios.

\end{cvparagraph}

\newpage

\null
\cvsection{Código: Detalle de contribuciones}

\begin{cventries}

%---------------------------------------------------------
  \cventry
    {Librerías y herramientas} % Job title
    {Contribuciones a proyectos Open Source} % Organization
    {} % Location
    {} % Date(s)
    {
    \begin{cvitems} % Description(s) of tasks/responsibilities
      \item {\textbf{\href{https://github.com/polirritmico/ebook-scriptum}{polirritmico/ebook-scriptum}}\\ Arquitectura modular para transformar textos simples o de lenguaje de marcado (ebooks, html, xml, xhtml, markdown, etc.) a través de LLM, TTS u otras herramientas de conversión (ffmpeg, lame, id tags, etc.).}
      \item {\textbf{\href{https://github.com/polirritmico/bakupipe}{polirritmico/bakupipe}}\\ Concepto de prueba para implementar un sistema de pipeline a través de scripts configurables con el entorno de desarrollo Godot. Aplica builds, tests unitarios, de aceptación, performance, etc., realiza un deploy a Github y a un repositorio de artefactos (Drive, GCP, etc.); y genera los logs.}
      \item {\textbf{\href{https://github.com/polirritmico/Bakumapu-docs}{polirritmico/Bakumapu-docs}}\\ Documentación html generada automáticamente a partir del documento LaTeX.}
      \item {\textbf{\href{https://github.com/polirritmico/ansible-bootstrap}{polirritmico/ansible-bootstrap}\href{https://github.com/polirritmico/ansible}{https://github.com/polirritmico/ansible}}\\ Configuración de entornos de desarrollo en Linux a través de técnicas de bootstraping con scripts Bash y playbooks Ansible (Ubuntu, Fedora, Gentoo, Arch).}
    \end{cvitems}
    }

%---------------------------------------------------------
  \cventry
    {Plasma y Qt (C++)} % Job title
    {KDE} % Organization
    {} % Location
    {} % Date(s)
    {
    \begin{cvitems} % Description(s) of tasks/responsibilities
      \item {\textbf{\href{https://store.kde.org/p/2075700}{https://store.kde.org}}\\ Build script para modificar el tema de cursores original de KDE.}
      \item {\textbf{\href{https://invent.kde.org/multimedia/elisa/-/merge_requests/478\#note_741609}{kde.org/multimedia/elisa}}\\ Implementar QStandard paths al reproductor de música Elisa (reproductor por defecto de KDE).}
      \item {\textbf{\href{https://bugs.kde.org/buglist.cgi?bug_status=UNCONFIRMED&bug_status=CONFIRMED&bug_status=ASSIGNED&bug_status=REOPENED&email1=ebray187&emailassigned_to1=1&emailcc1=1&emaillongdesc1=1&emailreporter1=1&emailtype1=substring&order=Importance&query_format=advanced}{bugs.kde.org}}\\ Varios reportes de errores.}
    \end{cvitems}
    }

%---------------------------------------------------------
  \cventry
    {Distribución Linux avanzada basada en paquetes fuentes y rolling release} % Job title
    {Gentoo Linux} % Organization
    {} % Location
    {} % Date(s)
    {
    \begin{cvitems} % Description(s) of tasks/responsibilities
      \item {\textbf{\href{https://github.com/polirritmico/gentoo-localrepo}{polirritmico/gentoo-localrepo}}\\ Repositorio con ebuilds personalizados}
      \item {\textbf{\href{https://github.com/polirritmico/gentoo-patches}{polirritmico/gentoo-patches}}\\ Colección de parches personalizados a diferentes aplicaciones. Por ejemplo Dolphin (Personalizar aspectos visuales como ancho de barras, espaciados, etc.), Hydrogen (soporte XDG), Aegisub (suporte XDG, utf-8), Firefox (Personalizar algunos maps), Ghostty (bugfixes no integrados al upstream), neovim (ajustes de treesitter), etc.}
      \item {\textbf{\href{https://forums.gentoo.org/search.php?search_author=ebray187}{forums.gentoo.org}}\\ Guías, tutoriales, documentación y soporte a usuarios.}
    \end{cvitems}
    }

%---------------------------------------------------------
  \cventry
    {Fork de Vim. Editor con capacidades de IDE altamente configurable (Lua)} % Job title
    {Neovim} % Organization
    {} % Location
    {} % Date(s)
    {
    \begin{cvitems} % Description(s) of tasks/responsibilities
      %\item {\href{https://github.com/nvimdev/dashboard-nvim/pulls?q=involves%3Apolirritmico+}{https://github.com/nvimdev/dashboard-nvim/pulls?q\=involves\%3Apolirritmico\+}:\\ Lista con interacciones en proyectos (PR, merges, issues, etc.).}
      \item {\textbf{\href{https://github.com/polirritmico/telescope-lazy-plugins.nvim}{polirritmico/telescope-lazy-plugins.nvim}}\\ Este plugin hace un parsing sobre la configuración del usuario, sigue importaciones de módulos y lo integra en un fuzzy finder.}
      \item {\textbf{\href{https://github.com/polirritmico/monokai-nightasty.nvim}{polirritmico/monokai-nightasty.nvim}}\\ Una implementación del tema Monokai en lua para neovim. Proporciona utilidades como conmutador de temas light/dark y utiliza un sistema de caché sencillo para optimizar los tiempos de carga. Para validar la caché comprueba la configuración del usuario, último commit hash y compara algunos objetos serializados.}
      \item {\textbf{\href{https://github.com/polirritmico/simple-boolean-toggle.nvim}{polirritmico/simple-boolean-toggle.nvim}}\\ Conmutador booleano customizable}
      \item
          {\textbf{\href{https://github.com/polirritmico/manual-tag-closer.nvim}{polirritmico/manual-tag-closer.nvim}}\\
            A través de tressiter (syntax node parser) generar cierres de tags,
            como ``end", ``done", ``\}", ``fi", ``\</tag>", etc.}
      \item {\textbf{\href{https://polirritmico.github.io/}{polirritmico.github.io}}\\ Blog en Ruby, Jekyll}
      \item {\textbf{\href{https://www.reddit.com/user/ebray187/}{reddit.com/user/ebray187}}\\ Documentación y soporte a usuarios.}
    \end{cvitems}
    }

\end{cventries}

